\documentclass[../../main.tex]{subfiles}
\begin{document}
% SOURCE p=235; printed=220; object=ch07-ex12; number=12; region=middle-lower
\begin{exercise}\label{exercise:ch07-uncountable-lines}
令 $X=\mathbb R\times\mathbb R_d$，其中 $\mathbb R_d$ 表示赋予离散拓扑的实数集。若 $f$ 为 $X$ 上的函数，令 $f^y(x)=f(x,y)$；若 $E\subseteq X$，令 $E^y=\{x:(x,y)\in E\}$。
\begin{enumerate}[label=\alph*.,ref=(\alph*)]
\item\label{item:ch07-lines-cc} $f\in C_c(X)$ 当且仅当所有 $f^y\in C_c(\mathbb R)$，且只有有限多个 $y$ 使 $f^y\ne0$。
\item\label{item:ch07-lines-infinite-mass} 在 $C_c(X)$ 上定义正线性泛函 $I(f)=\sum_{y\in\mathbb R}\int f(x,y)\,\mathrm dx$，令 $\mu$ 为其对应的 Radon 测度。则对任何使 $E^y\ne\varnothing$ 的 $y$ 有不可数多个的集合 $E$，有 $\mu(E)=\infty$。
\item\label{item:ch07-lines-failure-inner} 令 $E=\{0\}\times\mathbb R_d$。则 $\mu(E)=\infty$，但对每个紧集 $K\subseteq E$ 有 $\mu(K)=0$。
\end{enumerate}
\end{exercise}
% SUPPLEMENT: id=v14-ch07-ex12-solution; parent=ch07-ex12; kind=solution+proof-detail+motivation+hypotheses; origin=independent-derivation
\begin{solution}
\textbf{不可数条开闭直线（重构新增）。}每条 $L_y=\mathbb R\times\{y\}$ 都是开闭子空间，同胚于 $\mathbb R$；$X$ 是这些直线的拓扑不交并，因而 LCH。

(a) 若 $f\in C_c(X)$，其紧支集向离散空间 $\mathbb R_d$ 的连续投影是紧集，因而有限（用单点开覆盖可证）。所以只有有限条直线上 $f$ 非零；每条限制连续，支集是原支集与闭直线之交的闭子集，故紧。反过来，逐条连续等价于在这个拓扑不交并上连续；若只有有限个非零切片且它们分别紧支，则所有切片紧支集的有限并紧，包含全函数的支集。因此 $f\in C_c(X)$。

(b) 由 (a)，定义 $I(f)$ 的和实际只有有限项，每项积分有限；所以 $I$ 线性且正。记 $m$ 为实线 Lebesgue 测度。先证明，对每个开集 $U\subseteq X$，
\begin{align*}
\mu(U)=\sum_{y\in\mathbb R}m(U^y),
\end{align*}
这里非负的不可数和定义为所有有限子和的上确界。若 $f\prec U$，它只在有限条直线上非零，且 $0\leqslant f^y\leqslant\chi_{U^y}$，所以 $I(f)\leqslant\sum_y m(U^y)$；Riesz 的开集公式给出 $\mu(U)$ 不超过右端。反过来，对任意有限指标集 $J$，每个开集 $U^y$ 的 Lebesgue 测度都可由 $0\leqslant f^y\leqslant1$、紧支集位于 $U^y$ 内的连续函数积分取上确界得到。把这些有限个函数放在相应直线上，其他直线置零，得到 $f\prec U$。因此 $\mu(U)\geqslant\sum_{y\in J}m(U^y)$；再对 $J$ 取上确界，得公式。

若 $E$ 有不可数个非空切片，则任意开集 $U\supseteq E$ 的相应切片均为非空开集，因而有正的 Lebesgue 测度。不可数个正数的非负和必为无穷：若和有限，则每个不小于 $1/n$ 的项的指标集合有限，其可数并包含所有正项，与不可数矛盾。所以 $\mu(U)=\infty$。外正则性给出 Borel 集 $E$ 的 $\mu(E)=\infty$。原题的“任何集合”若不预设可测性，应严格表述为相应外测度 $\mu^*(E)=\infty$；上述对开外罩的论证同样证明这个更一般表述。

(c) $E$ 是闭集，各切片都是 $\{0\}$，故由 (b) 有 $\mu(E)=\infty$。其每个紧子集 $K$ 的第二坐标投影有限，且每条直线上至多一个点，所以 $K$ 有限。每个单点 $(0,y)$ 可由仅位于 $L_y$ 内、任意短的开区间覆盖；开集公式及外正则性给出该单点测度零。故 $\mu(K)=0$，这说明 $\mu$ 在 $E$ 上不内正则。

为后续题记录一个进一步结论：若 Borel 集 $H$ 只占可数条直线，则 $\mu(H)=\sum_y m(H^y)$。先对仅占一条直线的 Borel 集，利用该直线开闭，可以把开外罩限制在该直线内，外正则及开集公式便给出其测度等于切片的 Lebesgue 测度；再对可数条直线使用可数可加性即可。
\end{solution}
\begin{remark}
\textbf{动机与逐条件说明（重构新增）。}离散的第二坐标迫使每个紧测试只看有限条直线，而开集可以同时在不可数条直线上给出正长度。这使紧内逼近看不到轴 $E$ 的无穷质量。测度 $\mu$ 仍满足 Radon 的三项定义，因为它由 Riesz 定理构造；失败的是额外的\emph{所有 Borel 集}内正则性。由 (b)，每个有限测度 Borel 集只占可数条直线，故全空间不能 $\sigma$ 有限，这与\cref{corollary:ch07-sigma-finite-regular}一致。不可数和的解释不可省略，也不能把本测度直接误认为逐线 Lebesgue 测度求和；两者在轴上的值恰好不同。
\end{remark}
% END SUPPLEMENT: id=v14-ch07-ex12-solution

% SOURCE p=235; printed=220; object=ch07-ex13; number=13; region=lower
\begin{exercise}\label{exercise:ch07-zero-weight-failure}
在\cref{exercise:ch07-uncountable-lines}的背景下，令 $A=(\mathbb R\setminus\{0\})\times\mathbb R_d$，$\phi(x,y)=|x|$。则测度 $\mu_A(E)=\mu(A\cap E)$ 和 $\nu(E)=\int_E\phi\,\mathrm d\mu$ 都不是 Radon 测度。因此，\cref{exercise:ch07-restricted-radon}中 $\mu$ 为 $\sigma$ 有限、\cref{exercise:ch07-integrable-weight-radon}中 $\phi\in L^1(\mu)$、\cref{exercise:ch07-positive-continuous-weight}中 $\phi>0$ 这些假设都不能删除。
\end{exercise}
% SUPPLEMENT: id=v14-ch07-ex13-solution; parent=ch07-ex13; kind=solution+motivation+hypotheses; origin=independent-derivation
\begin{solution}
\textbf{零质量轴的所有开邻域却有无穷质量（重构新增）。}令 $E=\{0\}\times\mathbb R_d$。因为 $E\cap A=\varnothing$ 且 $\phi=0$ 于 $E$，有 $\mu_A(E)=\nu(E)=0$。任意开集 $U\supseteq E$ 在每条直线上都包含零点的开区间，故 $U\cap A$ 在不可数条直线上有非空切片。由上一题 (b)，$\mu_A(U)=\mu(U\cap A)=\infty$。

对加权测度，令 $H_n=U\cap\{|x|\geqslant1/n\}$。每条直线上的 $U$ 均含一个非零点，所以这些 $H_n$ 所占直线指标集合的可数并是整个 $\mathbb R$。若每个指标集合都可数，其并也可数，矛盾。故存在 $n$，使 $H_n$ 有不可数个非空切片，从而 $\mu(H_n)=\infty$。于是
\begin{align*}
\nu(U)\geqslant\int_{H_n}|x|\,\mathrm d\mu\geqslant n^{-1}\mu(H_n)=\infty.
\end{align*}
两测度在 $E$ 上为零，而所有包含它的开集测度均为无穷，故均不外正则，因此不是 Radon 测度。
\end{solution}
\begin{remark}
\textbf{三个条件的反例对应（重构新增）。}$A$ 甚至是开集，但 $\mu$ 不是 $\sigma$ 有限，故限制测度结论失效。密度 $\phi$ 连续非负，在每个紧集上有界，却不属于 $L^1(\mu)$，因为取 $U=X$ 的上述论证给出 $\int\phi=\infty$。它在轴上为零，因此也不满足严格正性。两测度都仍可数可加且在紧集上有限，失败的具体条件是外正则，不能笼统归因于“不是测度”。本例逐一保持各题其余假设，准确说明所指条件不能一般地删去。
\end{remark}
% END SUPPLEMENT: id=v14-ch07-ex13-solution

% SOURCE p=235; printed=220; object=ch07-ex14; number=14; region=lower
\begin{exercise}\label{exercise:ch07-semifinite-radon}
设 $\mu$ 为 $X$ 上的 Radon 测度，$\mu_0$ 为它的半有限部分（见\cref{exercise:ch01-semifinite-part}）。
\begin{enumerate}[label=\alph*.,ref=(\alph*)]
\item\label{item:ch07-semifinite-inner} $\mu_0$ 在所有 Borel 集上内正则。
\item\label{item:ch07-semifinite-finite-outer} $\mu_0$ 在每个满足 $\mu(E)<\infty$ 的 Borel 集 $E$ 上外正则。
\item\label{item:ch07-semifinite-cc-integral} 对每个 $f\in C_c(X)$，有 $\int f\,\mathrm d\mu=\int f\,\mathrm d\mu_0$。
\item\label{item:ch07-semifinite-lines-formula} 若 $\mu$ 为\cref{exercise:ch07-uncountable-lines}中的测度，$m$ 为实线 Lebesgue 测度，则对任意 Borel 集 $E$，有 $\mu_0(E)=\sum_{y\in\mathbb R}m(E^y)$。
\end{enumerate}
\end{exercise}
% SUPPLEMENT: id=v14-ch07-ex14-solution; parent=ch07-ex14; kind=solution+motivation+hypotheses; origin=independent-derivation
\begin{solution}
\textbf{紧集可见的半有限部分（重构新增）。}按定义，
\begin{align*}
\mu_0(E)=\sup\{\mu(B):B\in\mathscr B_X,\ B\subseteq E,\ \mu(B)<\infty\}.
\end{align*}
回顾这确为测度：对不交 $E_j$，每个有限测度 $B\subseteq\bigcup_jE_j$ 可按 $B\cap E_j$ 分解，给出 $\mu_0(\bigcup_jE_j)\leqslant\sum_j\mu_0(E_j)$；反向，对有限多个 $j$ 各取有限测度 $B_j\subseteq E_j$，它们的有限并仍有限测度，故 $\mu_0(\bigcup_jE_j)\geqslant\sum_{j=1}^N\mu(B_j)$。分别取上确界再令 $N\to\infty$，得到反向不等式。这也涵盖某上确界无穷的情形。

(a) 对每个有限测度 Borel 集 $B$，由\cref{proposition:ch07-sigma-finite-inner}，$\mu(B)=\sup_{K\subseteq B,\ K\text{ 紧}}\mu(K)$。所有紧集测度有限，故两次取上确界得到
\begin{align*}
\mu_0(E)=\sup_{K\subseteq E,\ K\text{ 紧}}\mu(K)
=\sup_{K\subseteq E,\ K\text{ 紧}}\mu_0(K),
\end{align*}
最后一步因为在有限测度集上 $\mu_0=\mu$。这就是内正则性。

(b) 若 $\mu(E)<\infty$，则 $\mu_0(E)=\mu(E)$，且 $\mu_0\leqslant\mu$，所以
\begin{align*}
\mu_0(E)\leqslant\inf_{U\supseteq E,\ U\text{ 开}}\mu_0(U)
\leqslant\inf_{U\supseteq E,\ U\text{ 开}}\mu(U)=\mu(E)=\mu_0(E).
\end{align*}
各项必须相等，证明外正则性。

(c) 记 $K=\operatorname{supp}(f)$。$\mu(K)<\infty$，两测度在 $K$ 的所有 Borel 子集上都相等。由简单函数积分及单调逼近，它们对支撑在 $K$ 内的非负函数积分相等，再作实虚正负分解便得 $f$ 的积分相等；两积分绝对收敛，因为 $f$ 有界且两者在 $K$ 上有限。

(d) 若 $B\subseteq E$ 且 $\mu(B)<\infty$，上一题背景中的切片结论说明 $B$ 只占可数条直线，因此
\begin{align*}
\mu(B)=\sum_y m(B^y)\leqslant\sum_y m(E^y).
\end{align*}
对 $B$ 取上确界，得 $\mu_0(E)$ 不超过右端。反向，固定有限指标集 $J\subseteq\mathbb R$ 及整数 $N>0$，取
\begin{align*}
B_{J,N}=E\cap\bigl([-N,N]\times J\bigr).
\end{align*}
它是 Borel 集，只占有限条直线，$\mu(B_{J,N})=\sum_{y\in J}m(E^y\cap[-N,N])\leqslant2N|J|<\infty$，故 $\mu_0(E)\geqslant\mu(B_{J,N})$。先令 $N\to\infty$，再对有限 $J$ 取上确界，即得反向不等式及所求公式。
\end{solution}
\begin{remark}
\textbf{动机与逐条件说明（重构新增）。}半有限部分收集所有有限质量子集，而 Radon 的有限测度集内正则性使它恰等于紧子集质量的上确界。因此它总内正则，却未必仍外正则。(b) 要求的是原测度 $\mu(E)<\infty$，不能偷换为 $\mu_0(E)<\infty$：在 (d) 的轴上 $\mu_0(E)=0$，所有开外罩的 $\mu_0$ 值却仍无穷。(c) 说明紧支连续测试区分不了这两测度，而表示定理的唯一性要求双方都 Radon；这里没有矛盾。(d) 用有限切片与有界截断证明，不依赖对非 $\sigma$ 有限乘积测度擅自应用 Fubini 定理。
\end{remark}
% END SUPPLEMENT: id=v14-ch07-ex14-solution
\subfile{section-02-ordinals.tex}
\end{document}
